\section{Conclusion and Discussion}
This paper presents a novel method designed for diffusion language models. The proposed \texttt{FreeCache} enables training-free caching with an off-the-shelf diffusion model. On top of that, \texttt{Guided Diffusion} facilitates the quality of diffusion by guiding the direct diffusion process with outstanding speedup. 
Different from prior works that relies on training or additional finetuning effort, the proposed method enables the end-to-end speedup by using the off-the-shelf diffusion model only. Compared to the baseline diffusion language models, the proposed method achieves an average of 12.14$\times$ speedup across different reasoning tasks with negligible accuracy degradation. More importantly, for the first time, our method enables diffusion model to achieve the comparable performance compared to the autoregressive language models while maintaining the similar accuracy. 

Despite these advantages, certain trade-offs remain: \texttt{FreeCache} trades accuracy and memory for speed, as its KV approximation introduces small but unavoidable errors and higher memory use. \texttt{Guided Diffusion} adds system complexity, requiring an extra AR guider, and the guider’s alignment does impact the effectiveness of the method.

% First, \textbf{FreeCache is an approximation that trades accuracy and memory for speed}. As a training-free method, its KV approximation introduces a small but inherent accuracy loss. While this loss can be reduced by increasing the block size or introducing intermediate recomputations, it is difficult to eliminate entirely without training-aware methods or fine-tuning. Additionally, caching the full KV projections increases the model's memory footprint, which could be a constraint on GPUs with limited VRAM.

% Second, \textbf{Guided Diffusion introduces system complexity}. Our two-model pipeline, while training-free, requires serving an additional lightweight AR guider, which complicates system scheduling and deployment. This complexity is increased when the DLM and guider use different tokenizers that must be bridged at inference time. The overall performance is also dependent on the availability and quality of a suitable AR guider model.

% Third, \textbf{our evaluation is based on a limited set of available large-scale DLMs}. While we demonstrate generality across both Dream-7B (adapted from an AR model) and LLaDA-8B (trained from scratch), the landscape of billion-parameter DLMs remains small. The effectiveness of our methods on future DLM architectures is an area for future work.

% Finally, we have not quantified the total resource overhead, such as the energy consumption of a two-model pipeline, or explored the sophisticated scheduling optimizations needed to fully realize these gains in a production environment.